% Exemplo: pôster (painel) de evento acadêmico, 90 x 120 cm, com a
% biblioteca poster do tcolorbox (ambiente tcbposter). Reproduz o pôster
% feito com beamerposter (poster-beamerposter.tex): faixa com o nome do
% evento, título, autor e seções em duas colunas. A classe é a article; o
% tcbposter divide o painel numa grade e cada seção é uma caixa
% (\posterbox) presa a essa grade. Toda a arte (faixas, ondas, ícone) é
% desenhada com TikZ e as figuras são geradas no próprio documento: não há
% imagens de terceiros.
% Compilar: latexmk -pdf poster-tcbposter.tex (biber para as referências)
% Testado com: tcolorbox 6.10, biblatex-abnt 3.4,
%   TeX Live 2024 ou mais recente
%%BEGIN:tcbposterpreambulo
\documentclass{article}
\usepackage[T1]{fontenc}
\usepackage[brazilian]{babel}
\usepackage{lmodern}    % serifa e matemática escaláveis (Figura 2)
% os símbolos grandes (somatório) ficam com 10 pt em qualquer tamanho no
% lmodern: aqui crescem com o texto
\DeclareFontFamily{OMX}{lmex}{}
\DeclareFontShape{OMX}{lmex}{m}{n}{<-> lmex10}{}
\usepackage{helvet}     % sem serifa, com as medidas da Arial
\renewcommand{\familydefault}{\sfdefault}
\usepackage[varqu,varl,scaled=0.95]{inconsolata}  % fonte de máquina do livro
% tamanho do painel (o evento define); sem margens: o tcbposter cuida delas
\usepackage[paperwidth=90cm, paperheight=120cm, margin=0cm]{geometry}
\usepackage[most]{tcolorbox}
\tcbuselibrary{poster}  % ambiente tcbposter e \posterbox
\usepackage{tikz}       % arte do painel e diagramas
\usetikzlibrary{positioning, calc}
\usepackage{pgfplots}   % gráfico de barras
\pgfplotsset{compat=1.18}
\usepackage{booktabs}
\usepackage{caption}    % \captionof: legendas fora de flutuantes
\usepackage{enumitem}   % espaçamento das listas
\usepackage{xurl}       % endereços longos quebram em qualquer ponto
\usepackage{ragged2e}   % texto justificado (\justifying)
\usepackage{csquotes}
% [S. l.] e [s. n.] saem quando faltam local ou editora (NBR 6023:2025)
\usepackage[backend=biber, style=abnt]{biblatex}
\addbibresource{poster-evento.bib}
% Referências: tamanho da letra, entrelinha e espaço entre uma referência
% e a seguinte. Para mudar a aparência da lista, altere só estes três
% valores.
% tamanho da letra (o texto do pôster tem 36 pt)
\newcommand{\refletra}{24pt}
% distância entre as linhas de uma mesma referência
\newcommand{\refentrelinha}{29pt}
% espaço entre as referências; a NBR 6023 pede uma linha em branco
% (\baselineskip), que num pôster ocupa muito espaço
\newcommand{\refespaco}{0pt}
% O biblatex compõe a lista com \bibfont (e não com o tamanho em vigor); o
% biblatex-abnt usa \normalfont no subtítulo, que aqui voltaria ao tamanho
% do texto: na lista, ele só tira o negrito e o itálico. O espaço entre as
% referências entra no \bibsetup: o estilo abnt o redefine ali, ao imprimir
% a lista, e por isso um \setlength comum no preâmbulo seria desfeito.
\renewcommand*{\bibfont}{\fontsize{\refletra}{\refentrelinha}\selectfont
  \renewcommand*{\normalfont}{\mdseries\upshape}}
\AtBeginDocument{\appto{\bibsetup}{\setlength{\bibitemsep}{\refespaco}}}
%%END:tcbposterpreambulo

%%BEGIN:tcbposterestilo
% Tamanhos de letra pedidos pelo evento: texto com 36 pt (legível a 1 m),
% títulos das seções com 40 pt, legendas com 30 pt
\renewcommand{\normalsize}{\fontsize{36}{44}\selectfont
  \abovedisplayskip=24pt plus 6pt minus 6pt
  \belowdisplayskip=\abovedisplayskip}
\renewcommand{\small}{\fontsize{30}{36}\selectfont}
\renewcommand{\large}{\fontsize{40}{48}\selectfont}
\pagestyle{empty}
\setlength{\JustifyingParindent}{0pt}  % \justifying sem recuo de parágrafo
\setlength{\parindent}{0pt}

% Cores do painel
\definecolor{eventoescuro}{HTML}{0B3C5D}
\definecolor{eventomedio}{HTML}{1D7A8C}
\definecolor{eventodestaque}{HTML}{E07A1F}
\definecolor{eventotexto}{HTML}{4A4A4A}

% Arte do painel, desenhada com TikZ (coordenadas em cm, a partir do
% canto inferior esquerdo da página): faixa do alto com o nome do evento,
% ícone de livro aberto e círculos; ondas no pé. Entra como underlay da
% caixa que cobre a página (opção coverage do tcbposter).
\newcommand{\nomeevento}{SIMPÓSIO DE EXEMPLO 2026}
\newcommand{\lemaevento}{Ensino, pesquisa e extensão}
\newcommand{\artedopainel}{%
  \begin{scope}[shift={(frame.south west)}, x=1cm, y=1cm]
    % faixa do alto, com a borda de baixo em onda
    \shade[left color=eventoescuro, right color=eventomedio]
      (0,120) -- (90,120) -- (90,106.5)
      .. controls (65,102.5) and (35,109) .. (0,104.5) -- cycle;
    % círculos translúcidos
    \foreach \x/\y/\r in {80/114/9, 69/119/5, 86/104/4, 8/118/4, 58/108/2.5}
      \fill[white, opacity=.07] (\x,\y) circle (\r);
    % ícone: livro aberto
    \begin{scope}[shift={(13.5,112.3)}, line width=4pt, white]
      \draw[fill=white, fill opacity=.15] (0,0)
        .. controls (-2.5,1) and (-5,1) .. (-7,0.2)
        -- (-7,-6.8) .. controls (-5,-6) and (-2.5,-6) .. (0,-7) -- cycle;
      \draw[fill=white, fill opacity=.15] (0,0)
        .. controls (2.5,1) and (5,1) .. (7,0.2)
        -- (7,-6.8) .. controls (5,-6) and (2.5,-6) .. (0,-7) -- cycle;
      \foreach \k in {1,2,3} {
        \draw[line width=2.5pt] (-5.8,-1.2*\k)
          .. controls (-4,-1.2*\k+.4) .. (-1.2,-1.2*\k-.2);
        \draw[line width=2.5pt] (5.8,-1.2*\k)
          .. controls (4,-1.2*\k+.4) .. (1.2,-1.2*\k-.2);
      }
    \end{scope}
    \node[anchor=west, text=white, font=\fontsize{80}{90}\selectfont\bfseries]
      at (25,113.2) {\nomeevento};
    \node[anchor=west, text=white, font=\fontsize{40}{48}\selectfont]
      at (25.3,108.6) {\lemaevento};
    % fio de destaque abaixo da faixa
    \draw[eventodestaque, line width=6pt] (6,101.5) -- (84,101.5);
    % ondas do pé
    \fill[eventomedio, opacity=.85] (0,0) -- (0,4.5)
      .. controls (25,8.5) and (55,1) .. (90,5.5) -- (90,0) -- cycle;
    \fill[eventodestaque, opacity=.9] (0,0) -- (0,2.2)
      .. controls (30,0.2) and (60,5.5) .. (90,2.5) -- (90,0) -- cycle;
  \end{scope}}

% Seção: caixa sem moldura (blankest) cujo título é desenhado como caixa
% arredondada com contorno na cor do evento; o texto vem justificado, com
% hifenização. \tituloalto centraliza o título na altura de 3,3 cm.
\newcommand{\tituloalto}{%
  \rule[\dimexpr-1.65cm+.5\fontcharht\font`T\relax]{0pt}{3.3cm}}
% secao=Título: opções da caixa (o título sai em maiúsculas)
\tcbset{secao/.style={blankest, enhanced, frame hidden, interior hidden,
    title={\MakeUppercase{#1}}, toptitle=0pt, bottomtitle=0pt,
    fonttitle=\large\bfseries, coltitle=eventoescuro, halign title=center,
    before title={\tituloalto}, title code={%
      \draw[eventomedio, line width=2.25pt, rounded corners=.65cm]
        ([shift={(1.125pt,1.125pt)}]title.south west) rectangle
        ([shift={(-1.125pt,-1.125pt)}]title.north east);},
    top=1.36cm, colupper=eventotexto, before upper={\justifying}}}

% Legendas numeradas acima da figura e da tabela (como pede a ABNT), com
% rótulo em negrito; \fonte escreve a fonte abaixo. Dentro das caixas não
% há flutuantes: o ambiente elemento centraliza e \captionof numera.
\captionsetup{font=small, labelfont=bf, labelsep=colon, skip=.15cm,
  position=above}
% [espaço]: distância do texto à legenda (opcional)
\newenvironment{elemento}[1][.56cm]{\par\vspace{#1}\centering}{\par}
% \fonte[espaço depois]{texto}
\newcommand{\fonte}[2][1.32cm]{\par\vspace{.4cm}%
  {\small Fonte: #2\par}\vspace{#1}}
% listas com marcador laranja, alinhadas à esquerda (itens curtos, como no
% pôster com beamer)
\setlist[itemize]{label=\textcolor{eventodestaque}{\textbullet},
  leftmargin=1.53cm, labelsep=.39cm, itemsep=3pt, topsep=0pt, parsep=0pt,
  before={\raggedright}}
% \cmd{nome}: comando do LaTeX destacado no código da Figura 2
\newcommand{\cmd}[1]{\textcolor{eventomedio}{\bfseries\textbackslash #1}}
% \marca{n}: marcador numerado, como os do livro (círculo laranja)
\newcommand{\marca}[1]{\tikz[baseline=(m.base)]\node[circle,
  fill=eventodestaque, text=white, inner sep=.12cm,
  font=\small\bfseries] (m) {#1};}
%%END:tcbposterestilo

%%BEGIN:tcbposterdados
% título em duas linhas (72 pt): \\ marca onde a linha quebra
\newcommand{\titulo}{LaTeX para documentos acadêmicos:\\ um livro e um guia
  aberto de exemplos testados}
\newcommand{\autor}{Rodrigo Smarzaro}
\newcommand{\instituicao}{Departamento de Exemplo, Universidade Exemplo
  (autor@exemplo.edu.br)}
% Apoio financeiro, centralizado no pé: troque as imagens de exemplo (pacote
% mwe) pelos logotipos dos financiadores; \apoiofalse retira a seção
\newif\ifapoio
\apoiotrue
%%END:tcbposterdados

\begin{document}
\normalsize
% Grade do painel: a página não tem margens; a caixa de cobertura
% (coverage) ocupa a página inteira, leva a arte no underlay e recua o
% conteúdo 6 cm dos lados, 19 cm do alto (faixa e fio) e 11,15 cm do pé
% (ondas). Sobra uma área de 78 x 89,85 cm, dividida em 2 colunas com
% 4,68 cm entre elas (como as colunas de 47% do beamer). As linhas (rows)
% da grade não são usadas: cada caixa fica abaixo (below) da anterior, com
% a altura do próprio conteúdo, e rowspacing dá o espaço entre elas; o
% apoio financeiro fica acima (above) do pé (bottom).
%%BEGIN:tcbpostergrade
\begin{tcbposter}[
  coverage={nobeforeafter, colback=white, left=6cm, right=6cm, top=19cm,
    bottom=11.15cm, underlay={\artedopainel}},
  poster={width=78cm, height=89.85cm, columns=2, rows=1, colspacing=4.68cm,
    rowspacing=1cm},
  boxes={blankest, colupper=eventotexto}]
%%END:tcbpostergrade

%%BEGIN:tcbpostertitulo
% \centering do LaTeX, não halign=center: o halign do tcolorbox fixa o
% espaço entre palavras no tamanho do texto (36 pt), e o título de 72 pt
% ficaria apertado
\posterbox[before upper={\centering}, bottom=1.6cm]{name=titulo, column=1,
  span=2, below=top}{%
  {\fontsize{72}{84}\selectfont\bfseries\color{eventoescuro}\titulo\par}
  \vspace{1.5cm}
  {\large\autor\par\instituicao\par}}
%%END:tcbpostertitulo

%%BEGIN:tcbposterintroducao
\posterbox[secao=Introdução]{name=intro, column=1, below=titulo}{%
  Trabalhos acadêmicos seguem normas detalhadas de formatação. O \LaTeX,
  criado por Lamport sobre o \TeX\ de Knuth, separa o conteúdo da forma:
  o autor escreve o texto e marca a estrutura, e o programa cuida da
  numeração, das referências cruzadas, das citações e da lista de
  referências \parencite{Lamport1994}. Mesmo assim, quem começa encontra
  material em português disperso, muitas vezes desatualizado, e exemplos
  que não compilam mais.}

\posterbox[secao=Objetivos]{name=objetivos, column=1, below=intro}{%
  \begin{itemize}
    \item Oferecer um livro em português sobre \LaTeX\ para trabalhos
      acadêmicos, artigos, apresentações e material didático.
    \item Manter um guia público, com licença aberta, em que todos os
      exemplos do livro compilam \parencite{latexacademico}.
    \item Mostrar o uso das normas da ABNT do início ao fim.
  \end{itemize}}
%%END:tcbposterintroducao

%%BEGIN:tcbpostermetodologia
\posterbox[secao=Metodologia]{name=metodologia, column=1, below=objetivos}{%
  Cada exemplo nasce no guia e é compilado por testes automáticos a cada
  versão, com o \TeX\ Live mais recente. O livro não copia o código:
  inclui o trecho do próprio arquivo do guia, delimitado por marcadores, e
  a figura do resultado gerada a partir dele (Figura~\ref{fig:fluxo}).
  \begin{elemento}
    \captionof{figure}{Do guia público ao livro}\label{fig:fluxo}
    \begin{tikzpicture}[caixa/.style={draw=eventomedio, line width=2pt,
        rounded corners=.5cm, minimum width=13cm, minimum height=4.5cm,
        align=center}]
      \node[caixa] (guia) {\textbf{Guia público}\\exemplos testados};
      \node[caixa, right=7cm of guia] (livro) {\textbf{Livro}\\texto e
        figuras};
      \draw[-latex, line width=4pt, eventodestaque] (guia) -- (livro)
        node[midway, above=.3cm, font=\small] {trechos e}
        node[midway, below=.3cm, font=\small] {resultados};
    \end{tikzpicture}
    \fonte{elaborado pelo autor.}
  \end{elemento}
  No livro, cada exemplo forma uma unidade com o código e o resultado lado
  a lado, e marcadores numerados ligam o texto às linhas do código: a
  seção \marca{1} é numerada sozinha, e a fórmula \marca{2} sai em
  destaque, centralizada (Figura~\ref{fig:unidade}).
  \begin{elemento}
    \captionof{figure}{Uma unidade Exemplo: código e resultado}%
    \label{fig:unidade}
    \begin{minipage}[c]{.56\linewidth}
      \begin{tikzpicture}
        \node[fill=black!5, draw=black!25, line width=1.5pt, inner sep=.8cm,
              font=\ttfamily\small, align=left] {%
          \begin{tabular}{@{}r@{\hspace{.5cm}}l@{}}
            \smash{\marca{1}} & \cmd{section}\{Introdução\} \\
                              & A média de \$n\$ valores é \\
            \smash{\marca{2}} & \cmd{[} \cmd{bar}\{x\} =
                                \cmd{frac}\{1\}\{n\} \\
                              & \hspace*{1.5em}\cmd{sum}\_\{i=1\}%
                                \textasciicircum\{n\} x\_i \cmd{]}
          \end{tabular}};
      \end{tikzpicture}
    \end{minipage}\hfill
    \begin{minipage}[c]{.4\linewidth}
      \begin{tikzpicture}
        \node[fill=white, draw=black!25, line width=1.5pt, inner sep=.8cm,
              align=left, text width=12.3cm, font=\rmfamily\small] {%
          {\bfseries 1\quad Introdução}\par\medskip
          A média de $n$ valores é
          \[ \bar{x} = \frac{1}{n}\sum_{i=1}^{n} x_i \]};
      \end{tikzpicture}
    \end{minipage}
    \fonte{elaborado pelo autor.}
  \end{elemento}}
%%END:tcbpostermetodologia

%%BEGIN:tcbposterresultados
\posterbox[secao=Resultados]{name=resultados, column=2, below=titulo}{%
  O livro e o guia cobrem os principais tipos de documento acadêmico
  (Tabela~\ref{tab:numeros}), e os exemplos de trabalhos acadêmicos
  seguem a NBR 14724 \parencite{NBR14724}, com a classe abnTeX2
  \parencite{abntex2}. A Figura~\ref{fig:temas} mostra os temas com mais
  arquivos de exemplo no guia.
  \begin{elemento}[1.2cm]
    % [t] e \vspace{0pt}: as duas legendas alinhadas pelo topo
    \begin{minipage}[t]{.46\linewidth}
      \vspace{0pt}\centering
      \captionof{table}{Livro e guia em números}\label{tab:numeros}
      \begin{tabular}{@{}lr@{}}
        \toprule
        \textbf{Item}          & \textbf{Total} \\
        \midrule
        Capítulos              & 17 \\
        Apêndices              & 4 \\
        Exemplos comentados    & 81 \\
        Trechos de código      & 164 \\
        Exemplos no guia       & 88 \\
        Receitas no guia       & 81 \\
        \bottomrule
      \end{tabular}
      \fonte[1cm]{elaborado pelo autor.}
    \end{minipage}\hfill
    \begin{minipage}[t]{.52\linewidth}
      \vspace{0pt}\centering
      \captionof{figure}{Exemplos por tema no guia}\label{fig:temas}
      \begin{tikzpicture}
        \begin{axis}[xbar, width=12cm, height=15.5cm, xmin=0, xmax=12,
            bar width=.9cm, symbolic y coords={Áreas, Figuras, Tabelas,
              Apresentações, Diversão, Referências, Básico, Código},
            ytick=data, axis x line=none, axis y line*=left,
            nodes near coords, every node near coord/.style={font=\small},
            % números fora do modo matemático: saem na fonte sem serifa
            /pgf/number format/assume math mode=true,
            tick label style={font=\small}]
          \addplot[fill=eventomedio!75, draw=none] coordinates {(5,Áreas)
            (6,Figuras) (6,Tabelas) (8,Apresentações) (8,Diversão)
            (8,Referências) (10,Básico) (10,Código)};
        \end{axis}
      \end{tikzpicture}
      \fonte[1cm]{elaborado pelo autor.}
    \end{minipage}
  \end{elemento}}
%%END:tcbposterresultados

%%BEGIN:tcbposterconclusoes
\posterbox[secao=Conclusões]{name=conclusoes, column=2, below=resultados}{%
  Com os exemplos testados a cada versão, o código do livro continua
  compilando quando os pacotes mudam \parencite{Mittelbach2023}. O guia
  está aberto a pedidos de exemplos e correções, e o leitor encontra nele
  o arquivo completo de cada exemplo, pronto para compilar.}

% referências no fim da segunda coluna (só o apoio financeiro ocupa a
% largura toda)
\posterbox[secao=Referências, top=1.06cm,
  before upper={\raggedright}]{name=referencias,
  column=2, below=conclusoes}{%
  \printbibliography[heading=none]}
%%END:tcbposterconclusoes

%%BEGIN:tcbposterapoio
% caixa com 45% da largura (35,1 cm de 78 cm), centralizada com xshift de
% (78 - 35,1)/2 = 21,45 cm a partir da coluna 1 e presa ao pé do painel
\ifapoio
\posterbox[secao=Apoio financeiro, top=.71cm, width=35.1cm,
  before upper={\centering}]{name=apoio, column=1, xshift=21.45cm,
  above=bottom}{%
  \includegraphics[height=3cm]{example-image-a}\hspace{2cm}%
  \includegraphics[height=3cm]{example-image-b}\hspace{2cm}%
  \includegraphics[height=3cm]{example-image-c}}
\fi
%%END:tcbposterapoio

\end{tcbposter}
\end{document}
